% =============================================================================
% TFM manuscript v02 — base: PDF 2026_FinalProject_Ales (1).pdf (17 pp)
% Orange (\revnew) = ediciones v02 posteriores al PDF
% Compile from docs/ (or set TEXINPUTS). Figures: see \graphicspath.
% =============================================================================
\documentclass[11pt]{article}

% -------------------------
% PACKAGES
% -------------------------
\usepackage[
a4paper,
left=2cm,
right=2cm,
top=2.3cm,
bottom=2.3cm
]{geometry}
\usepackage{times}
\usepackage{graphicx}
\usepackage{wrapfig}
\usepackage{float}
\usepackage{amsmath, amssymb}
\usepackage{booktabs}
\usepackage{caption}
\usepackage{subcaption}
\usepackage{hyperref}
\usepackage{setspace}
\usepackage{xcolor}
\usepackage{titlesec}
\usepackage{titling}
\usepackage{enumitem}
\usepackage{multicol}

\bibliographystyle{elsarticle-num}

% Revised text: \rev = integrated (no colour); \revnew = latest edits (highlighted)
\newcommand{\rev}[1]{#1}
\newcommand{\revnew}[1]{\textcolor{orange!75!black}{#1}}
\newenvironment{revnewblock}{%
  \begingroup\color{orange!75!black}%
}{%
  \endgroup%
}



\titlespacing*{\section}{0pt}{1.2ex}{0.6ex}
\setlength{\droptitle}{-2cm}

\pretitle{\begin{center}\LARGE\bfseries}
\posttitle{\par\end{center}\vspace{-0.8em}}
\preauthor{\begin{center}\normalsize}
\postauthor{\par\end{center}\vspace{-0.8em}}
\predate{}
\postdate{}

% Figures: Overleaf used imagen/; local repo has docs/figures_print and comparison figures.
\graphicspath{%
  {imagen/}%
  {figures_print/}%
  {../results/comparison/zone_2/all_models/figures/}%
  {../results/comparison/zone_2/all_models/figures/pm10graph/}%
  {../results/comparison/zone_2/all_models/figures/pred_vs_real/}%
}

\title{\textbf{Evaluation of Spatial Representation Strategies in Transformer Architectures for PM$_{10}$ Forecasting to Support Air Pollution Episode Management: A Case Study in Bajo Nervi?n}\\[0.4em]}

\setlength{\parindent}{0pt}
\setlength{\parskip}{0.35em}
\setlength{\textfloatsep}{7pt plus 2pt minus 2pt}
\setlength{\floatsep}{7pt plus 2pt minus 2pt}

\author{
\normalsize
Student: Ales Padr?\\
Supervisor: Javier del Ser\\
2026
}
\date{}

\begin{document}
\maketitle
\vspace{-0.5em}

% -------------------------
\section{Abstract}
\label{sec:abstract}

This work compares three Transformer-based spatial representation strategies for hourly PM$_{10}$ forecasting in the Basque Country (\textit{Bajo Nervi\'on}; $N{=}18$ stations) under a shared protocol: look-back $L{=}48$\,h (selected via ACF and confirmed by a $L\in\{24,48,72,96\}$ sensitivity analysis), horizon $H{=}24$\,h, chronological 70/10/20 split, and equal station-mean ranking by MA24 RMSE. Informer, GAT-Informer and AirFormer are evaluated across cumulative feature packs F1--F5. For the spatial contrast (H1--H2), dynamic spatial attention (AirFormer F3/F5) still beats the temporal baseline at most stations, whereas a fixed PM$_{10}$ graph (GAT-Informer) does not. With \textbf{XGBoost} as a strong non-Transformer tabular baseline, \textbf{two complementary rankings emerge}, depending on how Saharan intensity and episode memory are incorporated into F5. On the official F1--F5 ladder, \textbf{XGBoost F5} achieves the best global MA24 RMSE ($\approx$5.45; $\approx$+24\% skill versus MA24 persistence), ahead of AirFormer F5 ($\approx$5.65), Informer F1 ($\approx$6.06) and GAT-Informer F2 ($\approx$6.10), answering research questions H3--H4. \emph{Peak-oriented extensions} outside the ladder---AirFormer F3IL and XGBoost F5I---reach a similar global MA24 ($\approx$5.42), marginally improving on official XGBoost F5 while cutting high-PM$_{10}$ errors (80+~\textmu g\,m$^{-3}$ $\approx$94$\to$$\approx$80--83). A matched F3IL check (meteo+Saharan memory, no co-pollutants) favours AirFormer (5.42) over XGBoost (5.54), whereas F5 favours XGBoost (5.45) over AirFormer (5.65)---so each architecture peaks with a different feature recipe. Oracle/realistic routing analyses support deploying a single zone-level strong model rather than per-station selection. Yet even peak-oriented packs leave large absolute error in the rare 80+~\textmu g\,m$^{-3}$ regime (pooled hourly RMSE $\approx$80), limiting stand-alone operational deployment for extreme-episode decisions without expert oversight and forecast Saharan transport.
\revnew{A loss-function sensitivity analysis (MSE \emph{vs.}\ official MAE in AirFormer; Section~\ref{sec:evaluation-strategy}) yields only marginal gains in high-PM$_{10}$ bins and does not surpass XGBoost F5 or the F3IL/F5I packs; a single zone-level model with the appropriate feature pack suffices for decision support.}

\noindent\revnew{\textit{Leyenda v02:} texto en \textbf{negro} = PDF entregado (17\,pp); texto en \textbf{naranja} = ediciones v02 posteriores al PDF.}

%---------------------------------------------------------
\section{Introduction}
\label{sec:introduction}

Air pollution remains one of the most important environmental and public-health challenges worldwide. Short-term exposure to PM$_{10}$, PM$_{2.5}$, NO$_2$ and O$_3$ has been associated with increases in daily mortality and emergency hospital admissions due to respiratory and cardiovascular diseases, whereas long-term exposure contributes to the development of chronic respiratory and cardiovascular diseases, lung cancer and millions of premature deaths worldwide~\cite{Huang2024,EEA,RuizPaez2023}. Consequently, improving air quality has become a strategic priority for governments and international organizations, as reflected in the recommendations issued by the World Health Organization (WHO)~\cite{WHO}.

Traditionally, air quality management has relied on two complementary sources of information: monitoring networks, which provide continuous observations of atmospheric pollutant concentrations, and atmospheric dispersion or chemical transport models, which simulate pollutant emissions, transport, chemical transformation and removal. Together, these approaches constitute the basis of current air quality assessment and forecasting systems~\cite{WHO2023}. Nevertheless, accurate forecasting with sufficient lead time remains essential for reducing population exposure and supporting preventive actions~\cite{Gao2024}.

The adoption of the new \textit{European Air Quality Directive (EU) 2024/2881} reinforces the role of modelling and forecasting techniques in air quality assessment, environmental planning and public health protection~\cite{EUdir}. Likewise, the \textit{Spanish National Short-Term Action Plan} highlights the need for forecasting systems capable of anticipating pollution episodes and supporting timely mitigation measures~\cite{MITECO2021}. Accordingly, this work considers the 24-hour moving average of PM$_{10}$ as the operational forecasting target, using the corresponding MA24 error as the primary evaluation metric, complemented by hourly forecasting performance. Recent advances in artificial intelligence have opened new opportunities for air-quality forecasting. Deep learning models effectively capture nonlinear relationships and long-range temporal dependencies, but air pollution is inherently a spatio-temporal process in which concentrations depend not only on their temporal evolution but also on atmospheric transport, meteorology, topography, emission patterns and interactions between monitoring stations. Representing these spatial dependencies remains one of the main challenges in data-driven forecasting.

This challenge is particularly relevant in the \textit{Bajo Nervi\'on} area (Basque Country), where urban, industrial and port activities interact with the Bilbao estuary, sea-breeze circulations and complex terrain, generating highly dynamic atmospheric conditions. These characteristics make the area an appropriate case study for evaluating alternative spatial representation strategies in PM$_{10}$ forecasting.

Within this context, this work compares three Transformer-based forecasting approaches: a temporal Informer, the graph-based GAT-Informer and AirFormer, which incorporates structured spatial attention. Their performance is evaluated under a common experimental framework and compared with a competitive XGBoost baseline~\cite{Chen2016XGBoost}. The assessment focuses on MA24 forecasting accuracy, persistence-based skill and hourly performance to determine the operational value of each method for anticipating high-PM$_{10}$ episodes.

\paragraph{Research questions.}
The experimental design addresses four questions:
\begin{itemize}[leftmargin=1.4em,itemsep=0.25em]
\item \textbf{H1---Spatial representation versus temporal-only forecasting:} Do architectures with explicit spatial modeling mechanisms, namely GAT-Informer and AirFormer, improve PM$_{10}$ forecasts relative to a temporal Informer baseline under a common dataset, feature set and evaluation protocol?
\item \textbf{H2---Structured attention versus fixed graph connectivity:} Does AirFormer outperform GAT-Informer when representing inter-station dependencies in the \textit{Bajo Nervi\'on} area?
\item \textbf{H3---Baseline and episode management:} How do Transformers compare to the XGBoost baseline under equal station-mean MA24 ranking, and can they support episode management, including routing strategies and high-PM$_{10}$ peaks?
\item \textbf{H4---Vs persistence:} Do the evaluated models improve upon MA24 and last-observation persistence baselines?
\end{itemize}

\paragraph{Contributions.}
This work makes the following contributions:
\begin{itemize}[leftmargin=1.4em,itemsep=0.25em]
\item A common-protocol comparison of Informer, GAT-Informer and AirFormer for PM$_{10}$ forecasting in the \textit{Bajo Nervi\'on} area, using an input look-back of $L=48$ hours and a prediction horizon of $H=24$ hours.
\item Inclusion of XGBoost as a competitive non-deep-learning benchmark, trained under the same temporal splits, forecast horizons and feature packs. XGBoost is not treated as a fourth spatial strategy, but as a strong tabular baseline against which the value of Transformer-based spatial mechanisms can be assessed.
\item Evaluation based on MA24 RMSE, persistence skill, hourly error metrics, look-back sensitivity and performance during high-PM$_{10}$ episodes.
\item Assessment of deployment strategies by comparing oracle station-level selection with realistic zone-level model selection.
\end{itemize}

% -------------------------
\section{State of the art}
\label{sec:sota}
\subsection{Air Quality Forecasting Using Artificial Intelligence}

Air-quality forecasting has received increasing attention because of its relevance to environmental management, public-health protection and the implementation of short-term mitigation measures. Forecasting methods can broadly be grouped into physics-based atmospheric models and data-driven approaches.

Physics-based approaches, including atmospheric dispersion and chemical transport models, explicitly represent emissions, transport, chemical transformation and deposition processes. Their main advantage is physical interpretability and the possibility of simulating alternative emission or meteorological scenarios. However, these models require detailed and uncertain input information, substantial domain expertise and considerable computational resources~\cite{AQAT}. Data-driven methods have expanded rapidly with the availability of historical observations from automatic air-quality monitoring networks. Classical statistical models, such as autoregressive integrated moving average (ARIMA) models, have been widely used because of their simplicity and interpretability~\cite{Box2015}. Their performance may nevertheless be limited when pollutant dynamics are nonlinear, non-stationary and affected by complex interactions between emissions, meteorology and spatial transport.

Gradient-boosted decision trees, particularly XGBoost, are strong baselines for structured prediction problems, including environmental forecasting applications~\cite{Chen2016XGBoost,Wallek2024,Wang2023XGBoost}. When supplied with lagged pollutant concentrations, meteorological variables, calendar descriptors and station-related features, they can provide accurate forecasts at relatively low training cost. Although they do not explicitly model graph-based or attention-based spatial interactions, they provide a demanding non-deep-learning benchmark against which the added value of Transformer architectures can be assessed.

Deep-learning models have further extended the capacity of data-driven forecasting. Recurrent neural networks, especially Long Short-Term Memory (LSTM) and Gated Recurrent Unit (GRU) architectures, can capture nonlinear temporal dependencies and have been applied extensively to air-quality prediction~\cite{Hochreiter1997,Zhao2023,Qin2019}. More recently, Transformer architectures have emerged as a prominent approach for long-horizon time-series forecasting because attention mechanisms can capture relationships between distant time steps without the strictly sequential processing required by recurrent networks~\cite{Vaswani2017,wen2023transformerstimeseriessurvey}. For the 24-hour forecasting horizon considered in this work, with a 48-hour look-back window, Transformer models offer a suitable mechanism to learn dependencies across multiple diurnal cycles and meteorological regimes. Informer is a representative efficient Transformer architecture for long-sequence forecasting, using ProbSparse attention and sequence-distillation mechanisms to reduce computational complexity~\cite{Zhou2021Informer}. However, an Informer applied without an explicit spatial mechanism may not adequately represent the structure of interactions among monitoring stations.

\subsection{Spatial Representation for Air-Quality Forecasting}

Air pollution is an inherently spatio-temporal process. Measurements at a given station reflect local emission sources and meteorology, but they are also influenced by regional transport, topography, wind direction, atmospheric stability and pollutant concentrations at other locations. Consequently, the representation of spatial dependencies is a central modelling decision in air-quality forecasting.

A first family of approaches focuses primarily on temporal dynamics. In these models, each station is forecast independently, or information from multiple stations is included as multivariate input without an explicit spatial structure. This formulation provides an important baseline because it tests how much forecast skill can be obtained from local history and correlated covariates alone. Informer is used in this study as a representative temporal Transformer baseline~\cite{Zhou2021Informer,Zhou2023InformerEx}.

A second family represents the monitoring network as a graph. Stations are treated as nodes, while edges encode a predefined spatial relationship, such as geographic distance, $k$-nearest-neighbour connectivity, historical correlation or another similarity criterion. Graph Neural Networks aggregate information from connected stations, enabling the model to exploit structured inter-station dependence. Graph Attention Networks extend this principle by learning attention weights over the neighbours permitted by the graph connectivity~\cite{Velickovic2018}. Hybrid GAT-Informer models combine graph-based spatial aggregation with Transformer-based temporal forecasting and provide a natural framework for assessing the value of a fixed graph representation~\cite{YU2023110864}.

A third family uses attention mechanisms to model spatial relationships more flexibly. Rather than relying solely on a predefined adjacency matrix, these architectures learn spatial interaction patterns from data. AirFormer is a representative example: it combines temporal and spatial attention within a deterministic forecasting stage and incorporates a stochastic latent-variable stage to quantify uncertainty~\cite{liang2022airformer}. In its original evaluation, AirFormer was trained on four years of observations from 1,085 monitoring stations in China and reported improvements over comparison methods for 72-hour forecasts~\cite{liang2022airformer}. Its structured spatial-attention design is particularly relevant to settings where the relation between stations may vary with atmospheric conditions.

The distinction between fixed graph and attention-based methods should nevertheless be interpreted carefully. A graph-attention model can learn different weights among the neighbours allowed by its adjacency structure, whereas its potential interactions remain constrained by the initial graph. Conversely, attention-based approaches can model broader data-driven interactions, but learned attention weights should not be interpreted directly as physical transport pathways without additional atmospheric validation.

Despite substantial progress in spatio-temporal forecasting, published comparisons often differ simultaneously in geographical domain, monitoring-network density, input variables, prediction horizon, target pollutant, training/test splits and evaluation metrics~\cite{Agbehadji2024,Chen2024}. These differences make it difficult to attribute reported performance gains specifically to the spatial representation strategy. In particular, evidence remains limited on controlled comparisons of temporal-only Transformers, fixed-graph architectures and structured spatial-attention models for PM$_{10}$ forecasting in complex European urban-industrial environments.

This limitation motivates the present work. Informer, GAT-Informer and AirFormer are evaluated using the same study area, target variable, temporal splits, predictor variables, look-back window and forecasting horizon. XGBoost is trained under the same splits and feature packs as a competitive non-Transformer benchmark. The resulting experimental design aims to isolate, as far as possible, the contribution of alternative spatial representation strategies to PM$_{10}$ forecasting and to the operational estimation of the 24-hour moving average.

\subsection{PM$_{10}$ Forecasting and Episode-Oriented Evaluation}

Although PM$_{2.5}$ has received substantial attention in recent deep-learning air-quality forecasting studies, PM$_{10}$ presents specific modelling challenges. PM$_{10}$ may originate from combustion-related emissions, industrial activities, construction, road-dust resuspension, wind-blown dust, agriculture, wildfires and other natural or anthropogenic sources~\cite{Karagulian2015}. Road-dust resuspension is of particular relevance because it can contribute to the coarse-particle fraction and can vary with traffic activity, surface conditions and meteorology. The relative contribution of these processes may vary substantially between locations and atmospheric regimes. Consequently, findings obtained for PM$_{2.5}$ cannot be transferred directly to PM$_{10}$ without empirical validation. In particular, the relevance of coarse-particle sources and resuspension processes may increase the potential value of combining local temporal history with information from multiple monitoring stations and meteorological covariates.

For operational air-quality management, conventional global regression metrics may be insufficient. Recent PM$_{2.5}$ and PM$_{10}$ deep-learning studies commonly report RMSE and $R^2$ as primary evaluation metrics~\cite{Zhang2026X2AQFormer}. However, a model may exhibit low average error while systematically underestimating relatively infrequent high-concentration periods. Therefore, regression metrics should be complemented by high-concentration and event-oriented diagnostics. In particular, probability of detection (POD), false-alarm ratio (FAR), true skill score (TSS) and critical success index (CSI) have been applied to evaluate predictions of high PM$_{10}$ and PM$_{2.5}$ concentrations~\cite{Chae2021}.

The present work treats the forecasted 24-hour moving average of PM$_{10}$ (MA24) as its primary operational target. This choice is also consistent with the 24-hour averaging period used in the WHO Global Air Quality Guidelines~\cite{WHO} and Spanish National Short-Term Action Plan~\cite{MITECO2021}, which recommend a PM$_{10}$ concentration of 45~\textmu g\,m$^{-3}$ over 24 hours~\cite{WHO}. This health-based guideline should not, however, be interpreted as a substitute for the applicable legal thresholds or local episode-management protocols.

Accordingly, MA24 RMSE is used as the primary metric in this work. It is complemented by hourly MAE and RMSE, skill relative to persistence, and diagnostics focused on high-PM$_{10}$ observations. Where a specific episode threshold is defined, the evaluation can also include POD, FAR, TSS and CSI to assess the ability of the models to support episode-oriented decision-making.

% -------------------------
\section{Study Area and Dataset}
\label{sec:study}
\subsection{Study area}

The proposed methodology is evaluated using data from the official Air Quality Monitoring Network (AQMN) of the Basque Country~\cite{EuskadiAQMN} located in the \textit{Bajo Nervi\'on} (Figure~\ref{fig:bajonervion}) area, zone~2 according to regional methodological zoning.

\begin{figure}[H]
\centering
\includegraphics[width=0.60\textwidth]{bajonervion.jpg}
\caption{Study area and AQMN in the \textit{Bajo Nervi\'on}.}
\label{fig:bajonervion}
\end{figure}

Previous studies have shown that the study area constitutes a complex coastal atmospheric environment where pollutant transport is strongly influenced by the interaction between the Bilbao estuary, the surrounding orography and local meteorological processes~\cite{Oke1987,Whiteman2000,GonzalezAparicio2013}. This spatial variability cannot be explained solely by the temporal evolution at individual stations, motivating the explicit modelling of spatial dependencies (H2).

\subsection{Dataset}

The dataset was constructed by integrating hourly measurements from 18 monitoring stations collected within the study area. The data were retrieved from Open Data Euskadi, the open data portal of the Basque Government~\cite{OpenDataEuskadi}. All time series were synchronized using a common hourly temporal resolution, generating a unified spatio-temporal dataset suitable for model training.

The target variable considered throughout this work is PM$_{10}$, selected because of its environmental relevance and its role within the current European and Spanish regulatory framework for air quality management. In addition, several atmospheric pollutants, meteorological variables and temporal descriptors were incorporated as predictor variables in order to represent the main processes governing pollutant emission, transport and dispersion.

\begin{table}[H]
\centering
\caption{Variables included in the dataset.}
\label{tab:variables}
\renewcommand{\arraystretch}{1.15}
\begin{tabular}{ll}
\toprule
\textbf{Category} & \textbf{Variables} \\
\midrule
Target pollutant &
PM$_{10}$ \\
\addlinespace
Air pollutants &
PM$_{2.5}$, NO$_2$, NO, NO$_x$, SO$_2$, CO, O$_3$ \\
\addlinespace
Meteorological &
Wind speed, wind direction, temperature,\\
variables &
relative humidity, rainfall \\
\addlinespace
Temporal &
Hour of day, day of week, month \\
\addlinespace
Additional &
Saharan dust intrusion indicator \\
\bottomrule
\end{tabular}
\end{table}

Prior to model training, all variables were normalized and transformed into fixed-length temporal sequences using a sliding-window strategy, generating the input tensors required by the evaluated Transformer architectures.

% -------------------------
\section{Methodology}
\label{sec:methodology}
\subsection{Experimental Design}
The experimental design follows the same four aims introduced previously (H1--H4). Three Transformer-based architectures are compared under a common framework: Informer, GAT-Informer and AirFormer. To ensure a fair comparison, all architectures share the same study area, target pollutant, temporal resolution, input variables, prediction horizon, data partitioning strategy and evaluation metrics. Global ranking uses the equal station mean (1/18) of station-wise MA24 RMSE so that multi-node and per-station models are compared on the same footing. Therefore, the only experimental factor that changes among models is the strategy adopted to represent spatial relationships between monitoring stations. Figure~\ref{fig:experimental_workflow} summarizes the complete experimental workflow adopted in this work.

\begin{figure}[H]
\centering
\includegraphics[width=0.80\textwidth]{fig_experimental_workflow.png}
\caption{Experimental workflow.}
\label{fig:experimental_workflow}
\end{figure}

\subsection{Data Preparation and analysis}
Air quality and meteorological variables were downloaded through the official public API, synchronized to a common hourly temporal resolution and integrated into a unified spatio-temporal dataset. Since meteorological variables are not available at every monitoring station, missing observations were completed using a graph-based spatial assimilation strategy. A weighted graph was first constructed according to meteorological similarity among stations. Subsequently, missing meteorological observations were transferred from the most representative neighbouring stations before applying temporal interpolation to short gaps. This meteorological similarity graph is used for gap-filling only; it is not the adopted adjacency for GAT-Informer forecasting (Section~\ref{sec:methodology-gat}). An exploratory data analysis (EDA) was carried out to evaluate data availability, identify missing observations and analyse the relationships among variables. Based on this analysis, the common study period used for training is 2019-08-30 to 2025-12-29 (strict intersection of the 18 stations), hereafter referred to as late August 2019--December 2025.

\subsection{Feature Selection}
Feature selection followed a progressive cumulative ablation combining domain knowledge and exploratory data analysis~\cite{Zhang2026X2AQFormer}. Five nested feature configurations (F1--F5) were defined so that each level adds a variable block to the previous one, quantifying the contribution of calendar features, meteorology, co-pollutants and Saharan dust indicator. The Saharan indicator is the daily net African dust load (\textmu g\,m$^{-3}$) at the regional background site Valderejo, following the MITECO--CSIC procedure for natural PM$_{10}$ episodes~\cite{MITECO2013NaturalPM}.

\begin{table}[H]
\centering
\caption{Feature configurations evaluated in this work.}
\label{tab:features}
\begin{tabular}{ll}
\toprule
Configuration & Variables included\\
\midrule
F1 & PM$_{10}$\\
F2 & F1 + temporal variables\\
F3 & F2 + meteorology\\
F4 & F3 + co-pollutants\\
F5 & F4 + Saharan dust indicator\\
\bottomrule
\end{tabular}
\end{table}

Prior to defining these configurations, Autocorrelation Function (ACF) analysis was carried out to estimate the temporal memory of the PM$_{10}$ series. Based on the median autocorrelation across monitoring stations, an input sequence length of 48 hours was selected, while the prediction horizon was fixed at 24 hours. The 24\,h horizon matches the regulatory MA24 indicator used in the Spanish short-term action plan (H3). Importantly, despite this ACF-based choice of $L{=}48$\,h, a dedicated sensitivity analysis was performed by retraining the models with alternative look-back lengths $L\in\{24,48,72,96\}$ under the same splits, horizon and feature packs. Results (Table~\ref{tab:seqlen}) show that AirFormer remains essentially flat across these lengths, so $L{=}48$\,h is retained as the reference setting without loss of generality.

\subsection{Evaluated Architectures}
\label{sec:methodology-gat}
Three Transformer-based architectures were evaluated in this work. They differ exclusively in the way spatial information is represented, while sharing the same forecasting objective and experimental protocol. In addition, a per-station XGBoost model was trained on flattened lag windows with the same F1--F5 packs, chronological split and MA24 evaluation, providing a solid classical ML baseline. The starting implementations are the official repositories Informer~\cite{Zhou2021Informer,Zhou2023InformerEx}, AirFormer~\cite{liang2022airformer}, and GAT-Informer~\cite{YU2023110864}. The code was adapted to the experimental protocol of this work (pollutant, sequence length, prediction horizon, feature selection protocols, and metrics), while preserving the base architecture of each model.

\begin{table}[H]
\centering
\caption{Comparison of the evaluated architectures.}
\label{tab:architectures}
\renewcommand{\arraystretch}{1.15}
\begin{tabular}{p{3cm}p{3cm}p{3.2cm}p{2.5cm}}
\toprule
\textbf{Architecture} &
\textbf{Spatial representation} &
\textbf{Spatial relationships} &
\textbf{Prediction}\\
\midrule
Informer &
None &
Not explicitly modelled &
Per station \\
GAT-Informer &
Static graph &
Predefined graph &
Multi-station \\
AirFormer &
Dynamic attention &
Learned during training &
Multi-station \\
\bottomrule
\end{tabular}
\end{table}

Since all evaluated architectures provide forecasts for each monitoring station individually, the evaluation is first performed at station level. Global performance indicators are subsequently obtained by averaging the station-wise metrics, ensuring that each monitoring station contributes equally to the overall comparison regardless of the number of available observations.

\subsection{Training Protocol}
To ensure a fair comparison among the evaluated architectures, all experiments were conducted under a common training protocol. Prior to training, all predictor variables were normalized using statistics computed exclusively from the training dataset in order to avoid data leakage. The resulting normalization parameters were subsequently applied to the validation and test sets. The original PM$_{10}$ concentrations were recovered through the inverse transformation before computing the evaluation metrics. Input samples were generated using a sliding-window strategy with $L{=}48$\,h and $H{=}24$\,h. The complete dataset was divided chronologically into training (70\%), validation (10\%) and test (20\%) subsets. Each architecture was trained using the optimization strategy and hyperparameter configuration recommended in its original implementation, kept intentionally minimal to preserve the intended behaviour of each model. Due to the large number of experiments required, model training was performed on Tecnalia's High-Performance Computing (HPC) infrastructure (GPU/CPU). Table~\ref{tab:training} summarizes the common experimental settings.

\begin{table}[ht]
\centering
\caption{Common training settings adopted in the experimental evaluation.}
\label{tab:training}
\renewcommand{\arraystretch}{1.15}
\begin{tabular}{ll}
\toprule
\textbf{Parameter} & \textbf{Value} \\
\midrule
Target pollutant & PM$_{10}$ \\
Temporal resolution & 1 h \\
Input sequence length & 48 h \\
Prediction horizon & 24 h \\
Window generation & Sliding window \\
Train / Validation / Test & 70\% / 10\% / 20\% \\
Normalization & Training-set statistics \\
Model selection & Validation loss \\
Evaluation & Independent test set \\
Hardware & GPU \\
\bottomrule
\end{tabular}
\end{table}

\subsection{Evaluation Strategy}
\label{sec:evaluation-strategy}

The evaluated architectures were compared from two complementary perspectives. First, predictive performance was quantified using the Mean Absolute Error (MAE) and Root Mean Squared Error (RMSE) on hourly forecasts. Second, an operational assessment was carried out on the 24-hour mean of PM$_{10}$ derived from the predicted hourly concentrations---the primary ranking metric of this study. Metric definitions used throughout are: \emph{Hourly:} MAE/RMSE on the 24 forecast leads; \emph{MA24:} error between the mean of the 24 hourly forecasts and the mean of the 24 hourly observations for each window (primary for ranking); \emph{Skill:} relative improvement versus MA24 persistence (last observed 24\,h mean). This MA24-centred evaluation is the operational heart of H3. The same protocol was applied to the three architectures so that observed differences can be attributed to the spatial representation strategy. \revnew{As a supplementary \emph{loss-function sensitivity analysis}, AirFormer was additionally trained with mean squared error (MSE) instead of the official MAE of the original implementation, so that its training objective is comparable to the squared-error loss of the XGBoost baseline. Both loss variants use the same chronological splits, feature packs (F1--F5), early-stopping rule and test metrics; cross-model ranking remains based on MA24 RMSE. The outcome of this analysis is reported in Section~\ref{sec:results-peaks}.}

% -------------------------
\section{Results}
\label{sec:results}

All models were evaluated under a common protocol on the 18 PM$_{10}$ stations of the Basque Country using a 48\,h look-back window, a 24\,h prediction horizon, chronological train/validation/test splits and MA24 RMSE as the primary metric. Feature configurations F1--F5 progressively incorporate calendar variables, meteorology, co-pollutants and Saharan dust indicators. Informer2020, GAT-Informer, AirFormer and XGBoost were evaluated under identical conditions. Global rankings use the equal station mean (1/18), while persistence provides the H4 baseline.

Results are organised by research question. Table~\ref{tab:results-best} summarises the best official F1--F5 configuration for each model family and provides the common ranking used in H1--H2. Peak-oriented Saharan-intensity extensions further improve these results (Section~\ref{sec:results-peaks}).

\begin{table}[ht]
\centering
\caption{Best official F1--F5 configuration for each model family. Lower MA24 RMSE indicates better performance. Skill is computed relative to MA24 persistence. Peak-oriented configurations (F3IL/F5I) further improve these results (Table~\ref{tab:peak-ext}).}
\label{tab:results-best}
\small
\begin{tabular}{@{}llrrr@{}}
\toprule
Model & Best config & MA24 RMSE & Hourly RMSE & Skill (\%) \\
\midrule
\textbf{XGBoost} & F5 & \textbf{5.45} & 9.02 & \textbf{+24.4} \\
AirFormer & F5 & 5.65 & 9.08 & +21.6 \\
Informer2020 & F1 & 6.06 & 9.52 & +15.9 \\
GAT-Informer & F2 & 6.10 & 9.55 & +15.4 \\
Persistence (MA24) & --- & 7.20 & 10.37 & --- \\
Persistence (last) & --- & 8.79 & 11.51 & --- \\
\bottomrule
\end{tabular}
\end{table}

% ----- H1 -----
\subsection{H1 - Spatial versus temporal modelling}
\label{sec:results-h1}
\label{sec:results-spatial-map}

\textbf{Question:} Does incorporating spatial modelling (AirFormer, GAT-Informer) improve PM$_{10}$ forecasts relative to a purely temporal Informer?

Table~\ref{tab:results-best} shows the official F1--F5 ranking. AirFormer is the best Transformer, outperforming Informer, whereas GAT-Informer does not consistently improve over the temporal baseline. XGBoost provides the best overall result under the common protocol. \revnew{[Recorte v02: omitida Fig.~\ref{fig:ma24-by-config}; curva F1--F5 en Tab.~\ref{tab:global-mean-f}.]}

Global means alone do not show where spatial modelling helps. Figure~\ref{fig:spatial-benefit} maps, for each station, the MA24 RMSE difference $\Delta = \mathrm{RMSE}_{\mathrm{Informer}} - \mathrm{RMSE}_{\mathrm{model}}$ using the best F-config of each model (green $=$ model better than Informer; red $=$ Informer better).

\begin{figure}[H]
\centering
\includegraphics[width=0.98\textwidth]{fig_spatial_vs_temporal_map.png}
\caption{(Bajo Nervi\'on): station-wise MA24 RMSE benefit versus Informer (best F per model). Left: AirFormer; right: GAT-Informer. Green indicates lower error than Informer.}
\label{fig:spatial-benefit}
\end{figure}

AirFormer improves on Informer at most stations (largest gains at inland/mid-estuary sites such as San Juli\'an and Sangroniz; clear benefits around Muskiz--Abanto and the Bilbao core). A few coastal/northern sites (Algorta, San Miguel) are near-neutral. \revnew{[Recorte v02: omitidas Fig.~\ref{fig:heatmap}, Tab.~\ref{tab:best-station} y Fig.~\ref{fig:pred-overview}; AirFormer gana 17/18 estaciones (GAT-Informer solo en San Miguel F2).]}

With one model per station, AirFormer is the selected Transformer at 17/18 sites (mostly F5; F3 at a few stations); GAT-Informer F2 wins only at SAN MIGUEL. Predictions follow the observed MA24 baseline well, but even the local best Transformer under-predicts the sharpest high-PM$_{10}$ peaks---the residual H1 limitation that Section~\ref{sec:results-peaks} addresses with Saharan intensity/memory packs.

\textbf{Answer (H1):} Spatial modelling helps when it is dynamic: AirFormer F5/F3 (MA24 RMSE $\approx$5.65 / 5.70) beats Informer F1 ($\approx$6.06) globally and at most stations. A fixed-graph GAT-Informer does not systematically improve on the temporal baseline in \textit{Bajo Nervi\'on}. XGBoost F5 ($\approx$5.45) leads the official ladder but does not overturn this Transformer-level spatial--temporal contrast.

% ----- H2 -----
\subsection{H2 - Dynamic versus fixed spatial structure}
\label{sec:results-h2}
\label{sec:results-gat-variants}
\label{sec:results-ablation-f}

\textbf{Question:} Do dynamic spatial relations (AirFormer) better capture Bajo Nervi\'on dispersion than a fixed similarity graph (GAT-Informer)?

Beyond the H1 ranking, H2 examines how spatial structure is represented. Table~\ref{tab:global-mean-f} shows that AirFormer benefits from richer feature sets, reaching its best performance at F5, whereas GAT-Informer peaks at F2 and remains close to Informer. \revnew{[Recorte v02: omitida Tab.~\ref{tab:wins-f}; AirFormer domina desde F3 (16/18 en F3, 15/18 en F4--F5).]}

\begin{table}[ht]
\centering
\caption{Global mean MA24 RMSE by feature configuration and model ($n{=}18$ stations).}
\label{tab:global-mean-f}
\footnotesize
\begin{tabular}{@{}lrrr@{}}
\toprule
F & Informer & AirFormer & GAT-Informer \\
\midrule
F1 & 6.06 & 6.23 & 6.10 \\
F2 & 6.16 & 6.27 & 6.10 \\
F3 & 6.31 & 5.70 & 6.11 \\
F4 & 6.27 & 5.83 & 6.31 \\
F5 & 6.38 & 5.65 & 6.17 \\
\bottomrule
\end{tabular}
\end{table}

GAT-Informer represents the monitoring network using a fixed Pearson top-5 similarity graph, adopted because it provided the best forecasting performance and is more consistent with pollutant co-variability than a purely meteorological graph. In contrast, AirFormer learns spatial interactions directly from data through dynamic attention. Figure~\ref{fig:af-vs-gat-adj} illustrates this difference for three representative stations. Whereas GAT assigns fixed neighbour contributions throughout the forecasting period, AirFormer adapts the contribution of neighbouring stations according to the observed conditions, assigning greater weight to the target station while preserving a limited number of persistent neighbours.

Overall, AirFormer learns a self-centred, data-dependent spatial representation, whereas GAT relies on a fixed neighbour structure. This architectural difference explains the superior performance of AirFormer under the common protocol.

\begin{figure}[H]
\centering
\includegraphics[width=0.98\linewidth]{fig_af_vs_gat_spatial_adjacency.png}
\caption{Three example stations. Top: GAT-Informer---fixed PM$_{10}$ adjacency as ``\% that feeds the target'' (orange $=$ self). Bottom: AirFormer F5---projected attention to the same stations at 08:00, 12:00 and 18:00 (grouped bars). Three hours are shown for readability; the underlying attention is available for every hour of day.}
\label{fig:af-vs-gat-adj}
\end{figure}

\textbf{Answer (H2):} Dynamic spatial attention (AirFormer) outperforms the fixed PM$_{10}$ graph under the common protocol, both globally and across stations. AirFormer reaches its best official performance at F5 (MA24 RMSE 5.65), whereas GAT-Informer peaks at F2 and remains close to Informer. Figure~\ref{fig:af-vs-gat-adj} indicates that the advantage arises from learned spatial interactions rather than from the initial graph definition.

% ----- H3 -----
\subsection{H3 - Baseline ranking and episode management}
\label{sec:results-h3}
\label{sec:results-baseline}
\label{sec:results-management}
\label{sec:results-seqlen}

\textbf{Question:} Under equal station-mean MA24 ranking, how do the Transformers compare to the XGBoost baseline, and can the models support episode management under the shared protocol?

\paragraph{Transformers versus XGBoost.}
Table~\ref{tab:results-best} extends the comparison to the XGBoost baseline. Under the official F1--F5 ladder, XGBoost F5 achieves the lowest global MA24 RMSE (5.45), followed by AirFormer F5 (5.65), Informer F1 (6.06) and GAT-Informer F2 (6.10). AirFormer remains the strongest Transformer, whereas XGBoost provides the best overall performance.

Using the same F3IL feature configuration (meteorology + Saharan intensity and memory, without co-pollutants), AirFormer (5.42) outperforms XGBoost (5.54). Conversely, under the F5 configuration, XGBoost (5.45) outperforms AirFormer (5.65). These results indicate that each architecture benefits from a different feature combination rather than from a universally optimal predictor set.

\paragraph{Look-back robustness.}
Despite the ACF-based choice of $L{=}48$\,h, additional experiments with $L\in\{24,48,72,96\}$ confirmed that the main conclusions are insensitive to the selected look-back window (Table~\ref{tab:seqlen}). AirFormer remains stable across all lengths, XGBoost performs best at $L{=}48$, Informer favours shorter windows, and GAT-Informer deteriorates at $L{=}96$. Consequently, $L{=}48$\,h is retained as the reference setting.

\begin{table}[H]
\centering
\caption{Look-back sensitivity: best MA24 RMSE (equal station mean) per model and $L$. The best F-pack at each cell is shown in parentheses. Lower is better; the official protocol uses $L{=}48$.}
\label{tab:seqlen}
\footnotesize
\begin{tabular}{@{}lcccc@{}}
\toprule
Model & $L{=}24$ & $L{=}48$ & $L{=}72$ & $L{=}96$ \\
\midrule
AirFormer & 5.54 (F5) & 5.65 (F5) & 5.44 (F5) & 5.45 (F5) \\
XGBoost & 5.49 (F5) & \textbf{5.45} (F5) & 5.52 (F5) & 5.54 (F5) \\
Informer2020 & 5.95 (F1) & 6.06 (F1) & 6.03 (F1) & 6.06 (F1) \\
GAT-Informer & 6.23 (F3) & 6.10 (F2) & 6.09 (F1) & 7.41$^\dagger$ (F3) \\
\bottomrule
\end{tabular}
\end{table}

\paragraph{Oracle versus realistic routing.}
Oracle station-wise model selection improves MA24 RMSE by approximately 5\% relative to the best single global model. However, a realistic selector trained on the first half of the test period shows low agreement with the oracle and provides no meaningful improvement over a single zone-level model (Table~\ref{tab:oracle}). Consequently, a single global model is preferred for operational deployment. \revnew{[Recorte v02: fusionadas Tab.~\ref{tab:realistic} y Tab.~\ref{tab:oracle}; omitidas Fig.~\ref{fig:oracle} y Fig.~\ref{fig:realistic}.]}

\begin{table}[H]
\centering
\caption{\revnew{Oracle and realistic routing (MA24 RMSE). Oracle = best model$\times$config per station on full test; realistic = choose on 1st test half, evaluate on 2nd.}}
\label{tab:oracle}
\small
\begin{tabular}{@{}llrr@{}}
\toprule
Strategy & Config & MA24 RMSE & vs oracle (\%) \\
\midrule
Oracle (best per station, full test) & mixed & 5.19 & --- \\
XGBoost (single) & F5 & 5.45 & +5.0 \\
AirFormer (single) & F3IL & 5.42 & +4.2 \\
Realistic selector (eval half) & mixed & 5.24 & --- \\
XGBoost (single, best F on half) & F5I & 5.23 & --- \\
\bottomrule
\end{tabular}
\end{table}

\paragraph{High-PM$_{10}$ peaks and Saharan intensity extensions.}
\label{sec:results-peaks}
Episode management is driven by the upper tail. Stratifying hourly test error by observed PM$_{10}$ shows moderate errors in the 0--40\,\textmu g\,m$^{-3}$ bins, a sharp rise above 40\,\textmu g\,m$^{-3}$, and very large errors in the rare 80+ bin (Figure~\ref{fig:error-pm10}), often under Saharan-dust days. On the official F1--F5 ladder, F5 encodes continuous intensity and episode memory; XGBoost F5 (5.45) and AirFormer F5 (5.65) both improve on AirFormer F3 (5.70), but large errors persist in the 80+ bin.

\begin{figure}[H]
\centering
\includegraphics[width=0.98\linewidth]{fig_error_by_pm10_level.png}
\caption{Hourly error stratified by observed PM$_{10}$ level on the test set (headline models). Error concentrates in the high-concentration bins associated with Saharan episodes.}
\label{fig:error-pm10}
\end{figure}

To improve performance during extreme events, additional AirFormer configurations incorporating Saharan intensity (F3I), temporal memory (F3IL) and peak-weighted losses (F3IL\_pw) were evaluated under the same protocol. Table~\ref{tab:peak-ext} summarises global MA24. \revnew{[Recorte v02: omitida Tab.~\ref{tab:peak-slices}; bins 80+ ($\approx$94$\to$82) resumidos en texto.]}

\begin{table}[H]
\centering
\caption{Peak-oriented and official Saharan-intensity configs (equal station-mean MA24 RMSE). AirFormer F3I / F3IL / F3IL\_pw improve on official F3 ($\approx$5.70) and on official AirFormer F5 ($\approx$5.65$\to$5.42--5.48); XGBoost F5I matches $\approx$5.42. Matched F3IL (meteo+Sahara, no co-pollutants) is reported for both AirFormer and XGBoost.}
\label{tab:peak-ext}
\small
\begin{tabular}{@{}lrrr@{}}
\toprule
Config & MA24 RMSE & vs F3 & vs F3I \\
\midrule
F3 (official, AirFormer) & 5.70 & --- & --- \\
F5 (official, AirFormer) & 5.65 & $-0.05$ & --- \\
F3I ($+Sah\_int$, AirFormer) & 5.48 & $-0.23$ & --- \\
\textbf{F3IL} (AirFormer) & \textbf{5.42} & $\mathbf{-0.28}$ & $-0.06$ \\
F3I\_pw (aggressive peak loss) & 5.81 & $+0.10$ & $+0.33$ \\
\textbf{F3IL\_pw} (soft peak loss) & \textbf{5.43} & $\mathbf{-0.28}$ & $-0.05$ \\
F5 (official, XGBoost) & 5.45 & --- & --- \\
\textbf{XGBoost F5I} & \textbf{5.42} & --- & --- \\
XGBoost F3IL (meteo+Sahara) & 5.54 & --- & --- \\
\bottomrule
\end{tabular}
\end{table}

Peak-oriented configurations improve both global MA24 and extreme-event performance. AirFormer F3IL achieves the best Transformer MA24 (5.42), while XGBoost F5I reaches a comparable value; relative to F3, F3IL cuts 80+ hourly RMSE from $\approx$94 to $\approx$83. These improvements confirm the importance of Saharan intensity and episode-memory information for both Transformer and tree-based models~\cite{Wallek2024,Wang2023XGBoost}.

\revnew{The loss-function sensitivity analysis (Section~\ref{sec:evaluation-strategy}) shows that replacing MAE with MSE in AirFormer slightly improves high-PM$_{10}$ bins (40--80 and 80+) and the best global score at F5 (MA24 RMSE $\approx$5.58 \emph{vs.}\ $\approx$5.71 under MAE), but still does not beat XGBoost F5 ($\approx$5.45) or the peak-oriented packs F3IL/F5I ($\approx$5.42). From this analysis of alternative training losses, changing the objective alone does not justify added training or deployment complexity: a single zone-level model with the appropriate feature pack remains sufficient for decision support.}

\revnew{\paragraph{Lead-time error and short-horizon exceedance (H3).}
Although model ranking is based on MA24 RMSE, each 24\,h forecast is a vector of hourly leads. Figure~\ref{fig:lead-rmse-hit} decomposes this for the two official headline models (XGBoost F5 and AirFormer F5). Panel~(a) shows that hourly RMSE increases monotonically from $h{+}1$ to $h{+}24$ (from $\approx$5.3 to $\approx$10.2\,\textmu g\,m$^{-3}$ for XGBoost F5), so the first six hours are the most accurate part of the horizon. The dashed lines report MA24 RMSE ($\approx$5.45 / 5.65\,\textmu g\,m$^{-3}$): this metric is \emph{not} the average of hourly RMSEs, but the error on the 24-hour \emph{mean} concentration; averaging partially cancels hour-to-hour errors and therefore yields a substantially lower value than the long-lead hourly RMSE. Panel~(b) complements this with an exceedance-oriented view at 45\,\textmu g\,m$^{-3}$: the hit rate (POD) is highest at $h{+}1$ ($\approx$48\% XGBoost, $\approx$41\% AirFormer) and falls below 10\% beyond $h{+}6$, showing that even when hourly error is moderate, threshold detection degrades rapidly with lead time.}

\revnew{%
\begin{figure}[H]
\centering
\includegraphics[width=0.98\linewidth]{fig_lead_rmse_hit_headline.png}
\caption{Headline models (XGBoost F5, AirFormer F5): (a)~hourly RMSE by forecast lead (dashed lines = MA24 RMSE); (b)~hit rate (POD) for hourly PM$_{10}\geq 45$\,\textmu g\,m$^{-3}$ by lead (shaded: $h{+}1\ldots h{+}6$).}
\label{fig:lead-rmse-hit}
\end{figure}
}

\revnew{\paragraph{MA24 exceedance check ($\geq$ 45\,\textmu g\,m$^{-3}$).}
Short-term action plans in Spain often link preventive measures to sustained high PM$_{10}$ levels. Table~\ref{tab:ma24-45} extends the exceedance analysis to the operational MA24 indicator (mean of $h{+}1\ldots h{+}24$), using \textbf{XGBoost F5} and \textbf{AirFormer F3IL} (best peak-oriented Transformer). The MA24 hit rate ($\approx$35\%) is lower than the $h{+}1$ hourly POD in Figure~\ref{fig:lead-rmse-hit}b because MA24 aggregates all leads---including the less skilful long-range hours---into a single threshold test.}

\revnew{%
\begin{table}[H]
\centering
\caption{MA24 $\geq$ 45\,\textmu g\,m$^{-3}$ on the test set (zone~2, per station). \emph{Station rate:} fraction of station$\times$forecast windows with exceedance. \emph{Hit rate (POD):} among windows where the observed MA24 exceeds 45\,\textmu g\,m$^{-3}$, percentage also predicted above the threshold. \emph{Precision:} among predicted exceedances, percentage that are also observed above the threshold.}
\label{tab:ma24-45}
\small
\begin{tabular}{@{}lrr@{}}
\toprule
Metric & XGBoost F5 & AirFormer F3IL \\
\midrule
Observed station rate (\%) & 0.47 & 0.46 \\
Predicted station rate (\%) & 0.29 & 0.24 \\
Hit rate / POD (\%) & 34.8 & 35.4 \\
Precision (\%) & 54.4 & 66.1 \\
\bottomrule
\end{tabular}
\end{table}
}

\revnew{Exceedances are rare ($\approx$0.5\% of station-windows). Both models detect a similar share of observed MA24\,$\geq$\,45 events ($\approx$35\% hit rate), but under-predict the overall frequency (predicted rates 0.24--0.29\% \emph{vs.}\ $\approx$0.47\% observed). AirFormer F3IL is more conservative yet more precise when it does flag an exceedance (66\% \emph{vs.}\ 54\% for XGBoost F5). Together with Figure~\ref{fig:lead-rmse-hit}, these diagnostics show that short-lead hourly guidance is more skilful than long-lead or MA24-threshold alerts, but no metric alone is yet sufficient for automated action-plan activation.}

\textbf{Answer (H3):} Under the official protocol, XGBoost F5 provides the best overall performance, whereas AirFormer remains the strongest Transformer. Feature--architecture pairing proves more important than a single optimal feature set, and look-back sensitivity confirms the robustness of the 48\,h window. Peak-oriented configurations (AirFormer F3IL and XGBoost F5I) further improve MA24 and extreme-event performance. \revnew{By introducing continuous Saharan intensity and episode-memory features, the Transformer closes most of the gap to XGBoost on global MA24 (both $\approx$5.42; matched F3IL favours AirFormer, 5.42 \emph{vs.}\ 5.54), even though XGBoost still leads on the official F5 ladder ($\approx$0.20\,\textmu g\,m$^{-3}$).} Nevertheless, prediction errors during the most severe episodes remain too large for fully automated operational use.

% ----- H4 -----
\subsection{H4 - Skill versus persistence}
\label{sec:results-h4}

\textbf{Question:} Do learned models improve over MA24 / last-observation persistence baselines?

Table~\ref{tab:results-best} shows that all learned models outperform MA24 persistence. Under the official F1--F5 ladder, XGBoost F5 achieves the highest skill ($\approx$+24\%), followed by AirFormer F5 ($\approx$+22\%) and Informer/GAT ($\approx$+16\%). Peak-oriented configurations further increase skill to $\approx$+25\%. \revnew{[Recorte v02: omitida Fig.~\ref{fig:skill-pers}; skill en columna de Tab.~\ref{tab:results-best}.]}

\textbf{Answer (H4):} Yes. All learned models outperform MA24 persistence, with XGBoost F5 providing the highest official skill and AirFormer F3IL/XGBoost F5I further improving both persistence skill and high-PM$_{10}$ forecasting.

\section{Conclusions}
\label{sec:conclusions}

The experimental comparison identifies two complementary rankings. Under the official F1--F5 protocol, XGBoost F5 provides the best overall performance, whereas AirFormer is the strongest Transformer. When Saharan-intensity extensions are incorporated, AirFormer F3IL and XGBoost F5I achieve the best global MA24 performance, confirming that feature--architecture pairing is more important than a universally optimal feature configuration.

\begin{enumerate}[leftmargin=1.4em,itemsep=0.35em]
  \item \textbf{H1 (spatial vs temporal).} Dynamic spatial modelling improves PM$_{10}$ forecasting. AirFormer F5 consistently outperforms Informer, whereas GAT-Informer does not systematically improve over the temporal baseline.
  \item \textbf{H2 (dynamic vs fixed spatial structure).} Dynamic spatial attention captures inter-station relationships more effectively than a fixed PM$_{10}$ similarity graph. AirFormer consistently outperforms GAT-Informer under the common protocol, indicating that the performance gain is architectural rather than a consequence of the chosen adjacency graph.
  \item \textbf{H3 (baseline ranking and episode management).} XGBoost F5 provides the best official MA24 performance, whereas AirFormer is the strongest Transformer. \revnew{Peak-oriented Saharan packs (F3IL/F5I) narrow the Transformer--XGBoost gap and can match or slightly exceed the tree ensemble when intensity and episode memory are encoded explicitly (global MA24 $\approx$5.42).} \revnew{Complementary lead-time and exceedance checks (Figure~\ref{fig:lead-rmse-hit}, Table~\ref{tab:ma24-45}) show that hourly skill is highest in the first six hours (RMSE $\approx$5--8\,\textmu g\,m$^{-3}$) but threshold detection still captures only $\approx$35--48\% of high-PM$_{10}$ events, reinforcing that the models support---rather than replace---operational decision-making.} Nevertheless, prediction errors remain too large during the most severe events for fully automated operational use. The proposed models are therefore better suited as decision-support tools than as stand-alone episode management systems.
  \item \textbf{H4 (skill vs persistence).} All evaluated models outperform MA24 persistence. Peak-oriented configurations provide the highest skill, demonstrating that Saharan intensity and memory improve both persistence skill and high-PM$_{10}$ prediction.
\end{enumerate}

\revnew{The results do not require highly complex architectures for operational MA24 surveillance: XGBoost F5 matches or exceeds Transformers on the official ladder, and the loss-function sensitivity analysis (MSE \emph{vs.}\ MAE) does not change this ranking (Sections~\ref{sec:evaluation-strategy} and~\ref{sec:results-peaks}). Lead-time decomposition further shows that global MA24 scores integrate both accurate short-range hours and less skilful long-lead forecasts, while exceedance hit rates remain modest ($\approx$35--48\%) even for the best models---confirming their role as decision-support tools rather than autonomous alert systems.}

\subsection{Future work}
\label{sec:future-work}

\revnew{Future work follows directly from the operational caveat above: the present models are useful for decision support, but not yet for stand-alone episode management. The lines below prioritise closing that gap while preserving the reproducible protocol established in this thesis.}

\revnew{\textbf{Operational Saharan inputs.} The peak-oriented packs (F3IL/F5I) show that continuous Saharan intensity and episode memory materially improve both global MA24 and high-PM$_{10}$ bins. In deployment, however, these signals must be replaced by \emph{forecast} information. A natural next step is to ingest operational +24\,h dust products (e.g.\ CAMS mineral dust or MITECO Saharan alerts) and to cascade \emph{regime $\to$ concentration}: first classify the expected transport regime, then condition the PM$_{10}$ model on the predicted dust load rather than on observed intensity.}

\revnew{\textbf{Uncertainty and exceedance-oriented evaluation.} Ranking by equal-station-mean MA24 RMSE is appropriate for surveillance, but episode management requires metrics tied to regulatory thresholds. Future work should report probabilistic forecasts with calibration checks (e.g.\ Brier score) and deterministic exceedance diagnostics (POD, FAR, CSI) at 40, 50 and 80\,\textmu g\,m$^{-3}$, aligned with Basque alert levels. Conformal or ensemble-based uncertainty bands would help operators judge when a predicted exceedance is actionable rather than a false alarm.}

\revnew{\textbf{Regime-specific and mixture-of-experts models.} The complementary rankings of XGBoost (official F5) and AirFormer F3IL suggest that no single architecture dominates all concentration regimes. Mixture-of-experts or hard/soft regime routing---e.g.\ AirFormer under background conditions and gradient-boosted trees during Saharan episodes---could exploit the strengths identified here without increasing deployment complexity for every forecast cycle.}

\revnew{\textbf{Extended forecast horizon.} The official protocol uses $H{=}24$\,h; supplementary $H{=}48$\,h experiments (MA24\_D1 for the first day, MA24\_D2 for the second) indicate how skill degrades beyond the first daily mean. A full operational extension would re-tune look-back windows, loss weighting and feature packs for the two-day horizon and assess whether Saharan-memory features remain informative on MA24\_D2.}

\revnew{\textbf{Transferability.} The framework is demonstrated on zone~2 and PM$_{10}$. Applying the same feature-ladder, spatial-representation and evaluation protocol to other Basque air-quality zones and to additional pollutants (e.g.\ PM$_{2.5}$, NO$_2$) would test whether the AirFormer--XGBoost ranking and the value of Saharan features generalise beyond the Bajo Nervi\'on industrial--coastal setting.}

\revnew{\textbf{Adaptive spatial structure.} H2 showed that dynamic attention outperforms a fixed PM$_{10}$ similarity graph. Future work could learn time-varying adjacency (data-driven graph learning, attention over candidate edges) or fuse multiple graphs (distance, emissions, land use) instead of relying on a single pre-defined structure.}

\revnew{\textbf{Human--machine workflows.} Given the residual error in the 80+ bin, operational integration should treat model output as one input among others---meteorological briefings, CAMS dust maps and local expert judgement---rather than as an automated trigger. User-centred evaluation (forecast review sessions, post-event debriefs) would clarify how MA24 guidance is consumed in practice.}

\revnew{\textbf{Robustness.} Finally, sensitivity analyses should cover seasonal non-stationarity (winter stagnation \emph{vs.}\ summer dust), missing-station scenarios, and changes in the monitoring network. Stress-testing the F1--F5 ladder under realistic data gaps would strengthen confidence before any pilot deployment.}

\vspace{1em}
\noindent\revnew{\textbf{Data and code availability.} A reviewer-facing reproduction package (definitive Bajo Nervi\'on PM$_{10}$ dataset, analysis scripts, metric tables, manuscript figures and \texttt{REVIEWER.md} guide) is publicly available at \url{https://github.com/apadroh/2026_FPD_Ales_Padro_Bajo_Nervion}.}

\vspace{0.5em}
\noindent\textbf{Acknowledgements.} This work is supported by TECNALIA Research \& Innovation.

\noindent\textbf{Declaration of Generative AI.} During the preparation of this work, the authors used [GureGT-Qwen3] to translate the initial draft from Spanish and refine the grammatical flow of the manuscript. Additionally, generative AI tools, specifically [Cursor Version 3.13.10], were used under human supervision to accelerate the coding workflow.

\revnew{%
\hypertarget{fig:ma24-by-config}{}\label{fig:ma24-by-config}%
\hypertarget{fig:heatmap}{}\label{fig:heatmap}%
\hypertarget{tab:best-station}{}\label{tab:best-station}%
\hypertarget{fig:pred-overview}{}\label{fig:pred-overview}%
\hypertarget{tab:wins-f}{}\label{tab:wins-f}%
\hypertarget{tab:realistic}{}\label{tab:realistic}%
\hypertarget{fig:oracle}{}\label{fig:oracle}%
\hypertarget{fig:realistic}{}\label{fig:realistic}%
\hypertarget{tab:peak-slices}{}\label{tab:peak-slices}%
\hypertarget{fig:skill-pers}{}\label{fig:skill-pers}%
}

% -------------------------
\section*{References}
\label{sec:references}
\begin{multicols}{2}
\small
\bibliography{biblio,refs_metodologia}
\end{multicols}

\end{document}
